\begin{abstract}
\textbf{摘要。}
现代智能体并非独立运行的 LLM。它们运行在管理工具、上下文与控制流的\emph{智能体脚手架（agent harness）}之中，脚手架因此成为智能体的关键组成部分。我们最初的 Agent Lightning 工作提出了训练与智能体执行解耦的架构，通过 LLM 端点代理将任意智能体接入强化学习（RL）训练。近期的框架如 verl Uni-Agent、AReaL 2.0、slime v0.3.0 与 Polar 均沿用了这种基于代理的方式，使 RL 训练得以在脚手架中进行。本文将这一范式称为\emph{\add{\harl{}}}（Harnessed Agentic RL）：部署时脚手架直接参与模型后训练，从而缩小训练与实际使用之间的差距。

\vspace{\baselineskip}

我们发现\add{\harl{}}与传统智能体 RL 有本质区别，并带来一系列新挑战。传统智能体 RL 中，训练引擎拥有环境交互循环；而在\add{\harl{}}中，脚手架拥有这一循环，训练引擎只能观察到一串 LLM 请求-响应对。如何建模并将这些调用组装成训练样本仍是一个开放问题。通过仔细研究，我们识别出\harl{}的若干挑战，包括重分词、样本合并、优势计算、损失归一化与训练后端调度。我们发现，若处理不当，这些挑战会导致训练无效或不稳定。现有框架普遍对此语焉不详。本文首次对这些挑战给出系统性阐述。

\vspace{\baselineskip}

我们进一步提出 Agent Lightning v1.0，一个面向\add{\harl{}}的轻量级框架。我们将简洁作为第一原则，整个框架\textbf{仅用约 3500 行代码实现}。其紧凑设计支持任意智能体脚手架，并为研究上述挑战提供实用的试验平台。我们在通用指令遵循智能体、搜索智能体与编程智能体上验证了 Agent Lightning v1.0。对于编程智能体，我们发现现有 RL 框架支持有限：缺少数据与完整训练脚本，且依赖大规模计算资源。为填补这一空白，我们基于开源数据集与模型提供了完整的数据清洗管线与可复现的训练脚本。\textbf{仅使用 6K 训练样本与适量算力，RL 将 Qwen3.5-9B 在 SWE-bench Verified 上从 41.8\% 提升至 56.4\%，绝对提升 14.6\%。}我们发布完整工作流与脚本，以促进可复现的\add{\harl{}}研究。
\vspace{\baselineskip}

\end{abstract}
